% ---------------------------------------------------------------------------
%  QML Lab research project 04: Re-uploading Depth Study on MNIST
%  Build:  latexmk -pdf main.tex   (needs qmlproject.sty, harvardqml.bst, references.bib)
% ---------------------------------------------------------------------------
\documentclass[11pt,a4paper]{article}
\usepackage{qmlproject}
\qmlheader{Project 04: Re-uploading on MNIST}

\hypersetup{pdftitle={QML Lab project 04: Re-uploading Depth Study on MNIST}}

\begin{document}

\begin{center}
  {\small\textcolor{black!55}{QML Lab \textbullet\ research project 04}}\\[0.5em]
  {\LARGE\bfseries Re-uploading Depth Study on MNIST}\\[0.6em]
  {\large How do layers and qubits trade off in a small data re-uploading classifier for handwritten digits?}\\[0.6em]
  {\small Living document \textbullet\ last updated \today}
\end{center}

\begin{projectcard}
\small
\begin{tabularx}{\linewidth}{@{}>{\bfseries}l X >{\bfseries}l X@{}}
Quadrant  & \CQ\ (classical data, quantum processor) & Track & Variational \\
Difficulty & Easy & Tools & PennyLane, scikit-learn \\
Builds on & Module 5 of the QML Lab & Workspace & \path{qml_projects/reuploading-depth-mnist/} \\
\end{tabularx}
\end{projectcard}

\paragraph{Summary.}
A data re-uploading circuit encodes the input again before every trainable layer, which lets even a
few qubits express complicated functions. This project maps how test accuracy on pairs of
handwritten digits depends on the number of layers and qubits, and compares every configuration with
logistic regression and a small MLP on the same PCA features. It is the entry project of the lab and
already has a runnable experiment.

% ===========================================================================
\section{Background and key concepts}
% ===========================================================================
\subsection{Digits and PCA}
MNIST \citep{lecun1998} has $28\times28$ images of handwritten digits. The scikit-learn digits set,
used here first, is the smaller UCI collection of $8\times8$ images \citep{alpaydin1998}, with about
180 images per class. PCA reduces each image to $n$ components, one per qubit, scaled to $[0,\pi]$.

\subsection{Variational classifier}
As in Module 5, the model is the expectation value of $Z$ on the first qubit after a parameterised
circuit,
\begin{equation}
  f_\theta(x)=\braket{0|\,U(x,\theta)^\dagger Z_0\,U(x,\theta)\,|0}+b,\qquad
  U(x,\theta)=\prod_{\ell=1}^{L}W(\theta_\ell)\,S(x),
\end{equation}
where $S(x)$ is angle encoding and $W(\theta_\ell)$ a strongly entangling layer with $3n$ angles. The
classifier has $3nL+1$ trainable parameters \citep{schuld2020circuit}. Training minimises the
mean squared error to labels $\pm1$ with Adam \citep{kingma2015}.

\subsection{Why re-uploading helps}
\citet{schuld2021fourier} showed that such a model is a truncated Fourier series in the inputs,
\begin{equation}
  f_\theta(x)=\sum_{\omega\in\Omega}c_\omega(\theta)\,e^{i\omega\cdot x},
\end{equation}
and that each repetition of the encoding adds frequencies to $\Omega$. With one encoding the model is
a single sine per feature; with $L$ repetitions it reaches frequency $L$
\citep{perezsalinas2020}. Layers therefore buy expressivity, while qubits buy input features.

% ===========================================================================
\section{Research questions}
% ===========================================================================
The project is open-ended: these questions are the starting point, not the finish line. Add new
ones to the README and the journal as they come up.

\begin{rqlist}
  \item Where does adding layers stop improving test accuracy on 4 PCA features?
  \item For a fixed parameter budget, is it better to add qubits (features) or layers (depth)?
  \item How far is the best circuit from logistic regression and a small MLP on the same features?
\end{rqlist}

% ===========================================================================
\section{Method}
% ===========================================================================
\subsection{Data}
Digit pairs from the scikit-learn digits set: 3 vs 5 (default), 4 vs 9 and 1 vs 7. A stratified
70/30 train/test split per seed; PCA and scaling are fitted on the training part only. Later, full
MNIST reduced with PCA.

\subsection{Models}
The existing \texttt{baseline} experiment: $n$ qubits, $L$ re-uploading layers of
\texttt{AngleEmbedding} plus \texttt{StronglyEntanglingLayers}, MSE loss, Adam with learning rate 0.1,
mini-batches of 32, 15 epochs.

\subsection{Baselines}
Logistic regression on the same PCA features (already logged as \texttt{baseline\_accuracy});
an MLP and an RBF-SVM with matched parameter counts.

\subsection{Metrics}
\begin{itemize}
  \item Test accuracy (mean $\pm$ s.d. over 5 seeds) and final training loss.
  \item Number of trainable weights; training time per epoch.
\end{itemize}

% ===========================================================================
\section{Planned experiments}
% ===========================================================================
Each experiment becomes one \texttt{[[experiments]]} entry in the workspace, so its runs, parameters
and metrics are stored and compared in the portal.

\begin{table}[h]
\centering\small
\begin{tabularx}{\linewidth}{@{}l>{\raggedright}X>{\raggedright}X>{\raggedright\arraybackslash}X@{}}
\toprule
ID & Experiment & What varies & Main metrics \\
\midrule
  E1 & Depth sweep & $L\in\{1,2,4,6\}$, $n=4$ & test accuracy \\
  E2 & Qubit sweep & $n\in\{2,4,6,8\}$, $L=2$ & test accuracy \\
  E3 & Equal budget & $(n,L)$ pairs with $\approx$ same $3nL$ & accuracy per parameter \\
  E4 & Hard pairs & 3 vs 5, 4 vs 9, 1 vs 7 & gap to baselines \\
  E5 & Full MNIST & $28\times28$ images, PCA & accuracy, run time \\
\bottomrule
\end{tabularx}
\caption{Experiment plan. Results go into the journal and the results log (\cref{sec:log}).}
\end{table}

% ===========================================================================
\section{Running it in the QML Lab}
% ===========================================================================
The workspace lives in \path{qml_projects/reuploading-depth-mnist/}. Scripts import three helpers from
\path{qml_lab/plots.py}:
\begin{itemize}
  \item \texttt{params()} reads the parameters from the portal's run form (defaults when run by hand);
  \item \texttt{save()} stores figures with the run;
  \item \texttt{log\_metrics()} records numbers that appear in the run table.
\end{itemize}

\begin{lstlisting}
qml_projects/reuploading-depth-mnist/
  project.toml          metadata + experiments
  README.md             goal, open questions, ideas
  experiments/          runnable scripts
  docs/main.tex         this document
\end{lstlisting}

First experiment: This experiment already exists in the workspace and can be run from the project page.

\begin{lstlisting}[language={}]
[[experiments]]
slug = "baseline"
title = "Re-uploading classifier vs logistic regression"
script = "experiments/e01_baseline.py"
figure = "baseline.png"

[experiments.params]
digit_a = 3
digit_b = 5
n_qubits = 4
n_layers = 2
epochs = 15
seed = 0
\end{lstlisting}

\begin{lstlisting}[language=bash]
# from the portal folder: pick up changes to project.toml
python manage.py sync_projects
cd qml_projects/reuploading-depth-mnist
# run by hand with the default parameters
PYTHONPATH=../../qml_lab python experiments/e01_baseline.py
\end{lstlisting}

% ===========================================================================
\section{Quality bar}
% ===========================================================================
\begin{itemize}
  \item A classical baseline with a similar number of trainable parameters, tuned with the same care \citep{bowles2024}.
  \item At least five random seeds per configuration, reported as mean $\pm$ standard deviation.
  \item How results change with the number of qubits.
  \item At least one run under a realistic noise model.
  \item Fit PCA and scaling on training data only, so the test set stays unseen.
\end{itemize}

% ===========================================================================
\section{Open directions}
% ===========================================================================
\begin{itemize}
  \item Multi-class classification with one readout qubit per class.
  \item Amplitude encoding of the full $8\times8$ image on 6 qubits.
  \item Trainable input scaling, which changes the available Fourier frequencies.
  \item Repeat the best configuration under a noise model (Project 14).
\end{itemize}

% ===========================================================================
\section{Risks and limitations}
% ===========================================================================
\begin{itemize}
  \item Small test sets (about 100 images per pair) give noisy accuracies; always use several seeds.
  \item PCA features of easy pairs are almost linearly separable, so all models may hit a ceiling.
\end{itemize}

% ===========================================================================
\section{Results log}
\label{sec:log}
% ===========================================================================
Summarise findings here as they become stable; day-to-day notes belong in the portal journal.

\begin{table}[h]
\centering\small
\begin{tabularx}{\linewidth}{@{}l l >{\raggedright}X >{\raggedright\arraybackslash}X@{}}
\toprule
Date & Run & Finding & Next step \\
\midrule
 & & & \\
 & & & \\
 & & & \\
\bottomrule
\end{tabularx}
\end{table}

\bibliographystyle{harvardqml}
\bibliography{references}

\end{document}
